\section{The Tower}

There is one machine that builds the normed division algebras, and it works by doubling. Take an algebra with a conjugation, form ordered pairs of its elements, and multiply the pairs by a fixed rule, with a conjugation that flips the sign of the second slot. Each doubling produces an algebra of twice the dimension. Start from the one-dimensional line and double, and double, and double again.

You get the reals (here the integers), then the complex numbers in dimension two, then the quaternions in dimension four, then the octonions in dimension eight, then the sedenions in dimension sixteen. This is the Cayley-Dickson tower, and it is the only ladder there is.

Each rung pays a price. Every doubling quietly drops one law of arithmetic, always in the same order. Going from the reals to the complex numbers loses the order, the very i with square minus one that handed us the arrow of time. Going from the complex numbers to the quaternions loses commutativity, so now the order of multiplication matters. Going from the quaternions to the octonions loses associativity, so the grouping of a triple product starts to matter. And going from the octonions to the sedenions loses the norm itself, the multiplying length that made folding reversible.

That last loss is the cliff. When the norm dies, two nonzero things can multiply to zero, and you can write the pair down explicitly: in the sedenions, a definite sum of basis elements times another is exactly zero. So the tower of algebras that fold without losing information stops at the octonions, the last rung before the fall. The next section turns this ceiling and a matching floor into a single forced dimension.
